%=============================================================
% Paper 6: Hyatt & McEntarfer (2012), AER P&P
%=============================================================
\chapter{Hyatt \& McEntarfer (2012)}

\begin{paperbox}{Job-to-Job Flows in the Great Recession}
\textbf{Authors:} Henry Hyatt \& Erika McEntarfer (U.S.\ Census Bureau, CES)\\[3pt]
\textbf{Journal:} \textit{American Economic Review: Papers \& Proceedings}
102(3): 580--583\\[3pt]
\textbf{Year:} 2012 \quad\textbf{Field:} Labour economics; Macroeconomics\\[3pt]
\textbf{Classification:} Empirical (descriptive measurement)
\end{paperbox}

%------------------------------------------------------------
\section{Research Question}

How did flows of workers between jobs---both direct job-to-job transitions and
transitions involving nonemployment---change during the Great Recession, and
what happened to the earnings changes associated with those flows?

%------------------------------------------------------------
\section{Motivation}

Despite extensive documentation of aggregate unemployment dynamics during
the Great Recession, relatively little was known about the underlying worker-
level transitions---specifically, how the \emph{type} of job separation (direct
job-to-job versus separation through nonemployment) changed in severity and the
earnings consequences for each type. Standard labour-market surveys cannot trace
workers across employers. Two prior strands of work---Jacobson, Lalonde, and
Sullivan (1993) on earnings losses from displacement, and Topel and Ward (1992)
on early-career wage growth from job change---had documented that job change
matters for lifetime earnings but had not connected these to the cyclical
behaviour of transition types. Hyatt and McEntarfer construct the first
nationally-anchored longitudinal employer-employee measures of distinct
job-to-job flow categories to characterise labour market adjustment during
the Great Recession.

%------------------------------------------------------------
\section{Main Contribution}

\begin{enumerate}
  \item \textbf{Direct job-to-job (J2J) flows are strongly procyclical} and
    fell to a 12-year series low by early 2009, preceding the official
    recession start by 1--2 quarters. Separations to long nonemployment spells
    are countercyclical.
  \item \textbf{Earnings changes from job transitions are procyclical} at
    every transition type; in the Great Recession all transition types hit
    series lows and the earnings \emph{penalty} for nonemployment widened.
  \item \textbf{Housing bust case study:} Residential construction workers
    experienced a sharp J2J collapse (flows fell to 70\% of their
    2004--06 level), much higher rates of industry switching (two-thirds
    leaving construction), and median earnings losses of $-3.4\%$ for
    inter-industry movers versus small gains in prior periods.
\end{enumerate}

%------------------------------------------------------------
\section{Relationship to the Literature}

The paper sits at the intersection of (i) gross worker flows
(Blanchard--Diamond 1990; Davis--Haltiwanger 1992), (ii) J2J flow measurement
(Fallick--Fleishman 2004; Bjelland et al.\ 2011), and (iii) earnings dynamics
and job change (Jacobson, Lalonde, Sullivan 1993; Topel--Ward 1992). The
contribution is combining all three in a quarterly administrative data
framework that separates transition types by nonemployment duration.
It is the empirical predecessor to the Census Bureau's full J2J statistics
programme.

%------------------------------------------------------------
\section{Theoretical Framework and Economic Mechanism}

No formal model is presented. The organising framework rests on two forces:

\begin{itemize}
  \item \textbf{Procyclical voluntary quits:} In good times, workers
    make frequent J2J moves to capture wage gains up a job ladder;
    in recessions workers are reluctant to quit, collapsing J2J flows.
  \item \textbf{Earnings penalty amplified by spell duration:} Longer
    nonemployment spells generate larger earnings losses at reemployment,
    consistent with human capital depreciation, stigma, or mismatch.
\end{itemize}

%------------------------------------------------------------
\section{Data}

\begin{itemize}
  \item \textbf{Dataset:} LEHD (Longitudinal Employer-Household Dynamics)
    multi-state pilot database.
  \item \textbf{Frame:} 9 anchor states (CA, FL, GA, IL, KS, MI, NV, NC, ND);
    national job histories constructed for these workers covering flows to/from
    $>$40 states.
  \item \textbf{Period:} 1998--2010 (quarterly).
  \item \textbf{Flow types:} (1) J2J same quarter; (2) J2J adjacent quarter;
    (3) nonemployment 1 quarter; (4) nonemployment $>$1 quarter;
    (5) dominant separations (total).
  \item \textbf{Limitation:} Cannot distinguish unemployed from out of
    labour force; nonemployment durations only approximately measured
    in quarterly earnings data.
\end{itemize}

%------------------------------------------------------------
\section{Empirical Strategy}

Pure descriptive measurement:
\begin{enumerate}
  \item Classify separators by the duration of the intervening nonemployment
    spell between primary jobs.
  \item Compute seasonally adjusted quarterly time series for each flow type
    (Figure 1) and median earnings changes by transition type (Figure 2).
  \item Case study of residential building construction (NAICS 2361): compute
    frequency, industry destination, and earnings changes for J2J flows
    originating in construction across three periods (2001--03, 2004--06,
    2007--09).
\end{enumerate}

%------------------------------------------------------------
\section{Identification Strategy}

\textbf{No causal identification.} Pure descriptive measurement. The
comparison across periods for residential construction (Table~1) is a
descriptive before/after comparison, not an event study with a control group.

%------------------------------------------------------------
\section{Main Results}

\subsection*{Aggregate J2J Flows (Figure~1)}

\begin{itemize}
  \item Direct J2J flows (same and adjacent quarters) began falling in early
    2007---1--2 quarters \emph{before} the official recession start (NBER
    December 2007).
  \item Both types fell to a 12-year series low by early 2009.
  \item Separations to $>$1 quarter of nonemployment: slight countercyclical
    pattern; large spike in late 2008 driven almost entirely by spells lasting
    $\geq$1 year.
\end{itemize}

\subsection*{Earnings Changes from Job Transitions (Figure~2)}

\begin{center}
\small
\begin{tabular}{lcc}
\toprule
\textbf{Transition type} & \textbf{2006:Q2 (boom)} & \textbf{Recession} \\
\midrule
Direct J2J (same quarter) & $+9\%$ & Series low; penalty for NE widened \\
J2J (adjacent quarter) & $+3.8\%$ & Procyclical decline \\
Nonemployment 1 quarter & $0\%$ & Negative \\
Nonemployment 2--3 quarters & $-1.2\%$ & Larger losses than 2001 recession \\
\bottomrule
\end{tabular}
\end{center}

All transition types show procyclical comovement in earnings changes.

\subsection*{Residential Construction Case Study (Table~1)}

\begin{center}
\small
\begin{tabular}{lcc}
\toprule
\textbf{Measure} & \textbf{2004--06} & \textbf{2007--09} \\
\midrule
J2J flows (index) & 100 & $\approx$70 \\
Share staying in construction & $\sim$41\% & $\sim$37\% (two-thirds left) \\
Median wage change, within construction & $+5.2\%$ & $-0.6\%$ \\
Median wage change, cross-sector & $+2.7\%$ & $-3.4\%$ \\
Share flowing to Accommodation \& Food & 7.9\% & 10.0\% \\
\bottomrule
\end{tabular}
\end{center}

Accommodation \& Food Services movements associated with 13--21\% downward
earnings moves.

%------------------------------------------------------------
\section{Economic Magnitude and Interpretation}

The collapse in direct J2J flows begins before the NBER recession, suggesting
the labour market ``felt'' the downturn through job-switching channels first.
The 30\% fall in construction J2J flows implies that roughly one-third of the
job-changing activity that normally keeps construction workers on career
ladders stopped entirely. The earnings penalties---up to $-10.6\%$ for
retail trade movers; $-20\%$ for accommodation movers from construction---
illustrate how sectoral reallocation during the housing bust imposed large
permanent earnings losses.

%------------------------------------------------------------
\section{Mechanisms}

\begin{enumerate}
  \item \textbf{Procyclical voluntary quits:} J2J collapse reflects workers
    ceasing to quit voluntarily in recessions, not a change in the gross
    separation rate.
  \item \textbf{Earnings penalty amplified by nonemployment duration:}
    Consistent with human capital depreciation, stigma, or mismatch (workers
    forced into lower-quality matches under duress).
\end{enumerate}

%------------------------------------------------------------
\section{Robustness Checks}

Short P\&P paper; no explicit robustness section. The construction of J2J
flow categories follows established conventions from Bjelland et al.\ (2011)
and Fallick, Haltiwanger, and McEntarfer (2011). Restriction to
\emph{primary} (maximum earnings) jobs is motivated by prior evidence that
most meaningful J2J flows occur between primary jobs.

%------------------------------------------------------------
\section{Limitations}

\begin{enumerate}
  \item \textbf{P\&P format (4~pages):} No regression analysis, worker
    characteristic controls, or formal inference.
  \item \textbf{Nine-state frame:} Not nationally representative.
  \item \textbf{Cannot distinguish U from N:} LEHD shows only earnings.
  \item \textbf{Approximate nonemployment durations:} One full quarter
    of nonemployment = 3--8 months in reality.
  \item \textbf{No causality:} Housing bust case study is descriptive.
\end{enumerate}

%------------------------------------------------------------
\section{Policy Implications}

\begin{itemize}
  \item J2J flows begin declining before official NBER recession start
    $\Rightarrow$ fluidity measures should complement unemployment rate as
    leading labour market indicators.
  \item Large earnings losses for construction workers ($-3.4\%$ cross-sector;
    $-20\%$ to accommodation) support targeted retraining policies for
    displaced construction workers.
  \item J2J flows are the primary vehicle for early-career wage growth
    (Topel--Ward 1992); their cyclical collapse implies recession-scarring
    effects on youth wage trajectories.
\end{itemize}

%------------------------------------------------------------
\section{Overall Assessment}

A high-quality descriptive paper that opens up an important new data source
and establishes clean stylised facts about how the \emph{type} of worker
transition changes over the business cycle. Its strength is the data
construction and clarity of empirical patterns. Its weakness is the P\&P
format which prevents deeper analysis. The findings motivated the full Census
Bureau J2J statistics programme.

\textbf{Quality: Good for its format. Clean measurement, novel facts,
policy-relevant.}

%------------------------------------------------------------
\section{Relevance to My Research}

Directly relevant for Canadian RDC work. LEED or the Longitudinal Worker
File can be used to construct analogous J2J flow series for Canada. Testing
whether Canadian J2J flows are equally procyclical, and whether the Great
Recession or commodity bust produced similar collapse dynamics in resource
provinces, would be a natural paper.

%------------------------------------------------------------
\section{Potential Research Extensions}

\begin{itemize}
  \item Construct Canadian J2J flow series from LEED 2000--2024; compare
    procyclicality magnitude to U.S.\ and across provinces.
  \item Study J2J flows from resource-sector workers during commodity busts
    (2014--16 oil price crash; COVID-19).
\end{itemize}

%------------------------------------------------------------
\section{Research Gaps}

\begin{itemize}
  \item No formal distinction of supply-side (workers not quitting) vs.\
    demand-side (employers not hiring) mechanisms.
  \item No analysis of how J2J flows differ by worker age, education, or
    tenure---all relevant for career-ladder interpretations.
\end{itemize}

%------------------------------------------------------------
\section{Methodological Lessons}

\begin{enumerate}
  \item Restricting J2J flow measurement to \emph{primary} jobs (highest
    quarterly earnings) is crucial to avoid conflating secondary jobs with
    meaningful career transitions.
  \item Separating transitions by nonemployment duration (0, 1, $>$1 quarter)
    reveals very different cyclical patterns across types.
  \item Leading indicator property: J2J flows decline 1--2 quarters before the
    official NBER recession date.
\end{enumerate}

%------------------------------------------------------------
\section*{Five-Sentence Summary}

Hyatt and McEntarfer (2012) use a nine-state LEHD pilot database to construct
quarterly measures of four types of worker transitions---direct job-to-job,
short nonemployment, and long nonemployment---for the period 1998--2010. They
show that direct J2J flows are strongly procyclical, began declining in early
2007 before the official recession start, and fell to a 12-year low by early
2009, while separations to long nonemployment spells are countercyclical.
Earnings changes associated with all job transitions are procyclical, with
direct J2J movers earning $\approx$9\% more at reemployment in booms versus
near zero in recessions; the penalty for nonemployment widened in 2007--09.
A case study of residential construction reveals that J2J flows fell to
$\approx$70\% of their pre-recession level, two-thirds of movers left the
sector, and median earnings changes turned negative across all destinations.
The paper establishes that direct J2J flows---the primary vehicle for career
advancement---collapse disproportionately in recessions, with lasting
consequences for worker earnings trajectories.

\vspace{6pt}
\begin{quickref}
\begin{tabularx}{\linewidth}{@{}lX@{}}
\textbf{Key contribution} &
  First quarterly J2J flow series from LEHD separating transition types by
  nonemployment duration; shows J2J flows lead the recession; documents
  procyclical earnings changes with widening nonemployment penalty. \\[4pt]
\textbf{Key weakness} &
  Short P\&P format (4 pages); no statistical inference; nine-state frame
  not nationally representative; cannot distinguish U from N. \\[4pt]
\textbf{Most interesting gap} &
  The mechanism---workers not quitting vs.\ employers not hiring---is
  unresolved; identifying this distinction is crucial for policies targeting
  labour market fluidity. \\[4pt]
\textbf{Publishable extension} &
  Construct Canadian J2J flow series from LEED 2000--2024; study resource-
  sector workers during commodity busts (Alberta 2014--16) as a quasi-
  experiment for sectoral J2J collapse and earnings losses.
\end{tabularx}
\end{quickref}
